\begin{abstract}
Evaluating whether AI safety mitigations remain robust on sensitive scientific topics requires benchmarks that are themselves safe to release, resistant to contamination, and structured for auditing mitigation effectiveness rather than raw knowledge. We present the \textbf{Frontier Uplift Observatory}, a contamination-aware evaluation framework that addresses three gaps in current approaches: contamination blindness in public benchmarks, the absence of structured mitigation auditing, and the lack of release governance for evaluation artifacts. The framework uses a public/restricted split architecture with 24 public-safe evaluation items spanning 6 domain families and 6 reasoning types, scored across 5 safety-relevant metrics with a controlled error taxonomy. We evaluate six production models spanning all major frontier AI providers---GPT-4o, Claude Sonnet 4, Gemini 2.5 Pro, DeepSeek-V3, Llama-3.3-70B, and Qwen3-32B---under pre- and post-mitigation conditions, finding that safety system prompts improve all models but with substantial cross-model variance ($\Delta = +0.02$ to $+0.41$). Claude Sonnet 4 achieves the strongest post-mitigation profile (avg.\ 4.43/5), while models with lower safety baselines show the largest mitigation gains. A dual-scoring approach (pattern-based scorer + LLM-as-judge) validates scoring reliability with mean absolute divergence of 0.66 and 87\% error-tag agreement. The framework, evaluation data, and analysis pipeline are released as an open-source toolkit.
\end{abstract}
